\centering
\vspace{30mm}
\section*{\underline{Multhopp-Quadratur}}

\begin{tikzpicture}[remember picture, overlay]

% ---- ERSTE REIHE: Am Seitenursprung starten
\setlength{\rowheight}{50mm}

    % A1) Erste Kachel oben links
    \setlength{\cellwidth}{74mm}
    \Tile{A1}{([xshift=10mm,yshift=-66mm]current page.north west)}{\cellwidth}{\rowheight}
        {
            \vspace{-\TileSep}
            \titlerm{Indizierung}\\ \vspace{-4mm}
            \begin{align*}
                n=1\!:~&M\text{ Glieder des Fourier-Ansatzes} \\[-0.5ex]
                       &\text{(= Fourier-Koeffizienten)} \\[1.0ex]
                n=2\!:~&M\text{ Stützstellen, an denen die Funktions-} \\[-0.5ex]
                       &\text{werte }\gamma(\theta)\,\sin\!\left(n\,\theta\right)\text{ und die Zirkulations-} \\[-0.5ex]
                       &\text{stärken }\gamma(\theta)\text{ betrachtet werden} \\[1.0ex]
                n=3\!:~&M\text{ diskrete Stützstellen, an denen} \\[-0.5ex]
                       &\text{der induzierte sowie der geometrische} \\[-0.5ex]
                       &\text{Anstellwinkel ausgewertet werden}
            \end{align*}
        }{0}{1}{1}{0}

    % B1) Zweite Kachel rechts anbauen
    \setlength{\cellwidth}{58mm}
    \Tile{B1}{(A1.north east)}{\cellwidth}{\rowheight}
        {
            \vspace{-\TileSep}
            \titlerm{Diskretisierung}\\ \vspace{2mm}
            Äquiangulare Diskretisierung mit $M$\\
            Stützstellen mit $\mu = 1,\dots,M$:\\ \vspace{1mm}
            \begin{equation*}
                \theta_\mu = \frac{\mu\,\pi}{M + 1}
            \end{equation*}
            \begin{equation*}
                \eta_\mu   = \frac{2\,y_\mu}{b} = \cos\!\left(\theta_\mu\right) = \cos\!\left(\frac{\mu\,\pi}{M + 1}\right)
            \end{equation*}
        }{0}{1}{1}{0}

    % C1) Zweite Kachel rechts anbauen
    \setlength{\cellwidth}{58mm}
    \Tile{C1}{(B1.north east)}{\cellwidth}{\rowheight}
        {
            \vspace{-\TileSep}
            \titlerm{Fourieransatz}\\ \vspace{-2mm}
            \begin{align*}
                \gamma(\theta) &= 2\sum_{n=1}^M a_n\,\sin\!\left(n\,\theta\right) \\[3.0ex]
                           a_n &= \frac{1}{\pi}\,\int_0^\pi \gamma(\theta)\,\sin\!\left(n\,\theta\right)\,\text{d}\theta \\[0.5ex]
                               &\approx \sum_{\mu=1}^M \gamma_\mu\,\sin\!\left(n\,\theta_\mu\right)
            \end{align*}
        }{0}{0}{1}{0}

% ---- ZWEITE REIHE: Unten anbauen
\setlength{\rowheight}{39mm}

    % A2) Erste Kachel unten anbauen
    \setlength{\cellwidth}{60mm}
    \Tile{A2}{(A1.south west)}{\cellwidth}{\rowheight}
        {
            \titlerm{Beiwerte}\\ \vspace{-2mm}
            \begin{align*}
                C_\text{A}     &= \frac{\Lambda\,\pi}{M+1}\,\sum_{\mu=1}^M \gamma_\mu\,\sin\!\left(\theta_\mu\right) \\[1.0ex]
                C_\text{W\!,i} &= \frac{\Lambda\,\pi}{M+1}\,\sum_{\mu=1}^M \gamma_\mu\,\alpha_\mu\,\sin\!\left(\theta_\mu\right)
            \end{align*}
        }{0}{1}{1}{0}

    % B2) Zweite Kachel rechts anbauen
    \setlength{\cellwidth}{66mm}
    \Tile{B2}{(A2.north east)}{\cellwidth}{\rowheight}
        {
            \titlerm{Anstellwinkel}\\ \vspace{-4mm}
            \begin{align*}
                \alpha_{\text{i},\nu} &= b_{\nu\nu}\,\gamma_\nu - \sum_{\mu=1,\,\mu\neq\nu}^M b_{\nu\mu}\,\gamma_\mu \\
                                      &= \alpha_{g,\,\nu} - f_\nu\,\gamma_\nu \\
                \alpha_{\text{g},\nu} &= b_{\nu\nu} + f_\nu\,\gamma_\nu - \sum_{\mu=1,\,\mu\neq\nu}^M b_{\nu\mu}\,\gamma_\mu
            \end{align*}
        }{0}{1}{1}{0}

    % C2) Zweite Kachel rechts anbauen
    \setlength{\cellwidth}{64mm}
    \Tile{C2}{(B2.north east)}{\cellwidth}{\rowheight}
        {
            \titlerm{Gleichung}\\ \vspace{3mm}
            Hauptdiagonalen $b_\nu$ berücksichtigt\\
            und Nebendiagonal-Einträge mit $-1$\\
            multipliziert:
            \begin{equation*}
                \underline{\underline{b}}_\text{\,full} \cdot \underline{\gamma} = \underline{\alpha_\text{g}}
            \end{equation*}
        }{0}{0}{1}{0}

% ---- DRITTE REIHE: Unten anbauen
\setlength{\rowheight}{64mm}

    % A3) Erste Kachel unten anbauen
    \setlength{\cellwidth}{104mm}
    \Tile{A3}{(A2.south west)}{\cellwidth}{\rowheight}
        {
            \titlerm{Koeffizienten}\\ \vspace{-4mm}
            \begin{equation*}
                b_{\nu\nu} = \sum_{n=1}^M \frac{n}{M+1}\,\frac{\sin^2\!\left(n\,\theta_\nu\right)}{\sin\!\left(\theta_\nu\right)} = \frac{M+1}{4\,\sin\!\left(\theta_\nu\right)}
            \end{equation*}%
            \vspace{-4mm}
            \begin{align*}
                \text{Für }\left|\nu-\mu\right| = 1, 3, 5, ... :\quad b_{\nu\mu} &= -\sum_{n=1}^M \frac{n}{M+1}\,\frac{\sin\!\left(n\,\theta_\nu\right)\,\sin\!\left(n\,\theta_\mu\right)}{\sin\!\left(\theta_\nu\right)} \\
                &= \frac{\sin\!\left(\theta_\mu\right)}{(M+1)\left(\cos\!\left(\theta_\mu\right)-\cos\!\left(\theta_\nu\right)\right)^2} \\[2.0ex]
                \text{Für }\left|\nu-\mu\right| = 2, 4, 6, ... :\quad b_{\nu\mu} &= 0 \\[1.0ex]
                \text{Hauptdiagonale }b_\nu: \quad b_\nu &= b_{\nu\nu} + f_\nu = b_{\nu\nu} + \frac{2\,b}{\left(\frac{\text{d}C_\text{A}}{\text{d}\alpha}\right)_\text{2D}}\,t_\nu
            \end{align*}
        }{0}{1}{0}{0}

\end{tikzpicture}